% Section propre à la version complète : panorama des plateformes et des surveys
\section{Plateformes de la vision embarquée}
\label{sec:plateformes}

\subsection{Quatre familles de cibles}

Le corpus couvre quatre familles de plateformes, avec des outils et des métriques propres :
\begin{itemize}
  \item \textbf{SoC FPGA} (Zynq-7000, Zynq UltraScale+, Kria) : un accélérateur taillé pour
    le réseau, en précision arbitraire. Outils : Vitis~AI (moteur unique, INT8) ou flots
    dataflow (FINN, hls4ml, SATAY). Voir les sections~\ref{sec:architectures}
    et~\ref{sec:toolflows}.
  \item \textbf{GPU embarqués} (Jetson Nano, TX2, Xavier NX, Orin Nano et NX) : un
    accélérateur généraliste. Outil : TensorRT, en FP32, FP16 ou INT8. Selon la table de
    \citet[Table~4]{arxiv-2411.01431-efficient-deep-learning-infrastructures-embedded},
    une Jetson AGX Orin annonce 15 à 40\,W et 275\,TOPS (données constructeur).
  \item \textbf{NPU et VPU} en coprocesseur : Intel NCS2 et Myriad~2 (OpenVINO),
    Coral Edge TPU (INT8), Hailo-8 (HailoRT).
  \item \textbf{Microcontrôleurs}, avec ou sans accélérateur (STM32, GAP9, MAX78000) : moins
    de 1\,Mo de SRAM, un budget de quelques mW à quelques dizaines de mW.
    \citet[Table~4]{2025-el-zeinaty-designing-object-detection-models-tinyml} en recensent
    huit, dont quatre à accélérateur IA (STM32N6, GAP9, K210, MAX78000).
\end{itemize}
Les ordres de grandeur de puissance vont donc de quelques dizaines de milliwatts (VEGA :
62\,mW, \citealp[§V-D]{arxiv-2110.10486-tinyml-platform-device-continual-learning}) à
une quinzaine de watts (Jetson Orin NX : 10 à 14\,W de télémétrie,
\citealp[Table~4]{arxiv-2502.15737-performance-analysis-you-only-look}).

\subsection{Les surveys : l'embarqué sans FPGA, ou le FPGA sans comparaison}

Les huit surveys du corpus se répartissent en deux groupes qui ne se recoupent pas.

\paragraph{Surveys « embarqué ».}
\begin{itemize}
  \item \citet{arxiv-2411.01431-efficient-deep-learning-infrastructures-embedded} couvrent
    la conception manuelle, la NAS guidée par le matériel, la compression, l'apprentissage
    sur l'appareil et les LLM. Leur matériel se limite aux Jetson, TPU et NPU, sans aucun
    FPGA \src{Table~4}. Ils rappellent que les FLOPs ne prédisent pas la latence réelle
    \src{§2.3} et que l'entraînement embarqué est limité par la mémoire des activations,
    pas par les paramètres \src{§5.1}.
  \item \citet{2023-shuvo-efficient-acceleration-deep-learning-inference} couvrent toute la
    chaîne de l'inférence sur l'edge (Proc.\ IEEE). Le FPGA y occupe une sous-section
    \src{§XI-C1}, et la comparaison entre plateformes reste qualitative.
  \item \citet{2024-mittal-comprehensive-survey-deep-learning-lightweight} recense 26
    détecteurs légers. Le FPGA n'a qu'un paragraphe \src{§4.3.3}, et les FPS des
    Tables~5--6 ne sont associés à aucune plateforme, ce qui les rend inexploitables pour
    une comparaison matérielle.
  \item \citet{2025-el-zeinaty-designing-object-detection-models-tinyml} traitent la
    détection sur MCU. Le FPGA n'y est évoqué qu'en passant \src{§6.1.3, §8}.
\end{itemize}

\paragraph{Surveys FPGA.}
\citet{2025-hozhabr-survey-real-time-object-detection},
\citet{arxiv-2509.04153-real-time-fpga-cnns-detection}\pp{},
\citet{arxiv-1806.01683-accelerating-cnn-inference-fpgas-survey}\pp{} et
\citet{2018-venieris-toolflows-mapping-convolutional-neural-networks} ne comparent pas
quantitativement les FPGA aux Jetson ou aux MCU. Hozhabr et Giorgi le font de façon
qualitative \src{Table~5} ; les gains face au GPU qu'ils citent sont ceux des travaux
d'origine, souvent obtenus contre des GPU de bureau.

Aucun survey du corpus ne propose donc une synthèse quantitative FPGA / Jetson / MCU à NPU
sur des critères communs. La section~\ref{sec:comparaison} en esquisse une à partir de
\texttt{benchmarks.csv}, et montre pourquoi elle reste fragile.
